% Unit 108 English translation of chapter8.tex lines 282--453.
% Source: Methods of Algebra, Volume 2, commit
% 9a5803ff77dd3257484cb177f851a73770a59dd3 (CC BY 4.0).
% stable-unit-id: o014.aljabr2.chapter8.nerves-of-categories
% Coverage: complete Section 8.3, ``Example: Nerves of Categories'';
% stops immediately before Section 8.4 at line 454.

% segment-id: o014.li2.chapter8.nerves-of-categories.h001
\section{Example: Nerves of Categories}\label{sec:nerves}
\hypertarget{o014.aljabr2.chapter8.nerves-of-categories}{ }
% segment-id: o014.li2.chapter8.nerves-of-categories.p001
This section introduces the nerves of categories and, through them, the
classifying spaces of groups. Both are important examples of simplicial sets.

% segment-id: o014.li2.chapter8.nerves-of-categories.eg001
\begin{example}[Nerve of a category]\label{eg:nerve-cat}
	\index{nerve}
	\index[sym1]{NC@$\mathrm{N}\mathcal{C}$}
	Let $\mathcal{C}$ be a small category. Define the functor
	% segment-id: o014.li2.chapter8.nerves-of-categories.q001
	\[ \simpDelta^{\opp} \to \cate{Set}, \quad [n] \mapsto \left\{ \text{all functors } [n] \to \mathcal{C} \right\}. \]
	As a simplicial set it is denoted by $\mathrm{N}\mathcal{C}$ and called
	the \emph{nerve} of $\mathcal{C}$. Specifying an element of
	$\mathrm{N}\mathcal{C}_n$ is equivalent to specifying a functor
	$[n]\to\mathcal{C}$, that is, a chain of morphisms in $\mathcal{C}$
	% segment-id: o014.li2.chapter8.nerves-of-categories.q002
	\[ (f_1, \ldots, f_n): C_0 \xrightarrow{f_1} C_1 \xrightarrow{f_2} \cdots \xrightarrow{f_n} C_n. \]
	Thus $\mathrm{N}\mathcal{C}_0$ may be identified with
	$\Obj(\mathcal{C})$, and $\mathrm{N}\mathcal{C}_1$ with
	$\Mor(\mathcal{C})$.

	The face morphisms $d_i:\mathrm{N}\mathcal{C}_n\to
	\mathrm{N}\mathcal{C}_{n-1}$ and degeneracy morphisms
	$s_j:\mathrm{N}\mathcal{C}_n\to\mathrm{N}\mathcal{C}_{n+1}$ are given
	by
	% segment-id: o014.li2.chapter8.nerves-of-categories.q003
	\begin{align*}
		d_i(f_1, \ldots, f_n) & = \begin{cases}
			(f_2, \cdots, f_n), & i = 0 \\
			(\ldots, f_{i+1} f_i, \ldots), & 0 < i < n, \\
			(f_1, \ldots, f_{n-1}), & i = n,
		\end{cases} \\
		s_j(f_1, \ldots, f_n) & = \begin{cases}
			\left( \identity_{C_0}, f_1, \ldots, f_n \right), & j = 0 \\
			(\ldots, f_j, \identity_{C_j}, f_{j+1}, \ldots ), & 0 < j < n \\
			(f_1, \ldots, f_n, \identity_{C_n}), & j = n.
		\end{cases}
	\end{align*}
\end{example}

% segment-id: o014.li2.chapter8.nerves-of-categories.r001
\begin{remark}
	Unfolding the definitions shows that the simplicial sets
	$\mathrm{N}(\mathcal{C}^{\opp})$ and $\mathrm{N}\mathcal{C}$ are related
	by the order-reversing duality of
	Remark~\ref{rem:simpDelta-order-reversal}. The details are left to the
	reader.
\end{remark}

% segment-id: o014.li2.chapter8.nerves-of-categories.p002
The nerve realizes a category concretely through combinatorial/topological
data and contains all the information of the original category. This is the
content of Proposition~\ref{prop:nerve-ff}, which will soon be proved. We
first characterize which simplicial sets are nerves; the argument is
instructive.

% segment-id: o014.li2.chapter8.nerves-of-categories.pr001
\begin{proposition}[Characterization of nerves]\label{prop:inner-horn-filling}
	Let $X$ be a simplicial set. The following statements are equivalent.
	\begin{enumerate}[(i)]
		\item There is a small category $\mathcal{C}$ such that
		$\mathrm{N}\mathcal{C}\simeq X$.
		\item $X$ has the unique inner horn-filling property: for every integer
		$0<i<n$ and every morphism $\sigma':\Lambda^n_i\to X$ in
		$\cate{sSet}$ (such a $\Lambda^n_i$ is called an inner horn), there is
		a unique $\sigma:\Delta^n\to X$ extending $\sigma'$.
	\end{enumerate}
\end{proposition}
% segment-id: o014.li2.chapter8.nerves-of-categories.pf001
\begin{proof}
	First prove (i) $\implies$ (ii). Let $X=\mathrm{N}\mathcal{C}$ and
	$0<i<n$; we wish to extend $\sigma':\Lambda^n_i\to X$ to $\Delta^n$.
	For every $0\leq k\leq n$ (or $0<k\leq n$), the morphism
	$\Delta^0\xrightarrow{\{k\}}\Delta^n$ (or
	$\Delta^1\xrightarrow{\{k-1,k\}}\Delta^n$) factors through
	$\Lambda^n_i$. Denote its image under $\sigma'$ by
	$C_k\in X_0=\Obj(\mathcal{C})$ (or
	$[g_k:C_{k-1}\to C_k]\in X_1=\Mor(\mathcal{C})$). This gives an element
	of $X_n$,
	% segment-id: o014.li2.chapter8.nerves-of-categories.q004
	\[ C_0 \xrightarrow{g_1} C_1 \xrightarrow{g_2} \cdots \xrightarrow{g_n} C_n. \]
	Write $\sigma$ for the corresponding morphism $\Delta^n\to X$. The
	construction also shows that $\sigma$ is the only possible extension of
	$\sigma'$. Since
	$\Lambda^n_i=\bigcup_{j\neq i}\Image[\mathrm{d}^j:
	\Delta^{n-1}\to\Delta^n]$, it remains to verify
	% segment-id: o014.li2.chapter8.nerves-of-categories.q005
	\begin{equation}\label{eqn:inner-horn-filling-aux0}
		\sigma \circ \mathrm{d}^j = \sigma' \circ \mathrm{d}^j : \Delta^{n-1} \to X, \quad j \neq i .
	\end{equation}

	By the definition of the nerve, proving
	\eqref{eqn:inner-horn-filling-aux0} reduces to the following verification.
	For each $j\neq i$ and each pair of adjacent elements $h<k$ in the
	sequence $0,\ldots,\widehat{j},\ldots,n$ (where $\widehat{j}$ means that
	$j$ is omitted), the pullbacks of $\sigma$ and $\sigma'$ along
	$\Delta^1\xrightarrow{\{h,k\}}\Lambda^n_i\subset\Delta^n$ must agree.
	\begin{itemize}
		\item If $k=h+1$, this is guaranteed by the construction of $\sigma$.
		In particular, \eqref{eqn:inner-horn-filling-aux0} always holds when
		$j\in\{0,n\}$.
		\item Suppose $(h,k)=(j-1,j+1)$. If $n=2$, this case cannot occur. If
		$n>2$, then either $j-1>0$ or $j+1<n$. When $j-1>0$,
		$\{j-1,j+1\}$ factors as
		% segment-id: o014.li2.chapter8.nerves-of-categories.q006
		\[ \Delta^1 \xrightarrow{\{j-2, j\}} \Delta^{n-1} \xrightarrow{\mathrm{d}^0 = \{1, \ldots, n\}} \Lambda^n_i \subset \Delta^n. \]
		By the already known $j=0$ case of
		\eqref{eqn:inner-horn-filling-aux0}, the two pullbacks agree. Similarly,
		when $j+1<n$, the issue reduces to the known case $j=n$.
	\end{itemize}

	Now prove (ii) $\implies$ (i). Define a category $\mathcal{C}$ by
	$\Obj(\mathcal{C}):=X_0$ and, for all $C,C'\in X_0$,
	% segment-id: o014.li2.chapter8.nerves-of-categories.q007
	\[ \Hom_{\mathcal{C}}(C, C') := \left\{ f \in X_1: d_1(f) = C, \; d_0(f) = C' \right\}. \]
	Use the degeneracy map $s_0:X_0\to X_1$ to define the identity morphism
	$\identity_C$ as $s_0(C)$. We also depict a morphism $f$ as
	$0\xrightarrow{f}1$ to emphasize that it corresponds to a $1$-simplex
	$[1]=\{0,1\}\to X$.

	If $d_0(f)=d_1(g)$, define the composite by $gf:=d_1(\sigma)$, where
	$\sigma:\Delta^2\to X$ is the unique extension of
	% segment-id: o014.li2.chapter8.nerves-of-categories.q008
	\[\begin{tikzcd}[column sep=small]
		& 1 \arrow[rd, "g"] & \\
		0 \arrow[ru, "f"] & & 2
	\end{tikzcd}: \Lambda^2_1 \to X
	\quad \text{that is,}\quad d^0(\sigma) = g, \quad d^2(\sigma) = f. \]

	The laws $f\circ\identity_C=f$ and $\identity_{C'}\circ f=f$ are
	witnessed, respectively, by the following elements of $X_2$.
	% segment-id: o014.li2.chapter8.nerves-of-categories.q009
	\[ \begin{tikzcd}[row sep=large]
		& 1 \arrow[rd, bend left, "{\identity_{C'} = d_0 s_1(f)}"] & \\
		0 \arrow[rr, "{}" name=A, "{f = d_1 s_1(f)}"'] \arrow[ru, bend left, "{f = d_2 s_1(f)}"] & & 2 \arrow[phantom, from=1-2, to=A, "{s_1(f)}" description]
	\end{tikzcd} \quad \begin{tikzcd}[row sep=large]
		& 1 \arrow[rd, bend left, "{f = d_0 s_0(f)}"] & \\
		0 \arrow[rr, "{}" name=A, "{f = d_1 s_0(f)}"'] \arrow[ru, bend left, "{\identity_C = d_2 s_0(f)}"] & & 2 \arrow[phantom, from=1-2, to=A, "{s_0(f)}" description]
	\end{tikzcd}\]

	For associativity $h(gf)=(hg)f$, construct
	$\sigma':\Lambda^3_2\to X$ whose three faces are, respectively,
	% segment-id: o014.li2.chapter8.nerves-of-categories.q010
	\[\begin{tikzcd}[column sep=small]
		& 1 \arrow[rd, "g"] & \\
		0 \arrow[ru, "f"] \arrow[rr, "gf"'] & & 2
	\end{tikzcd} \quad \begin{tikzcd}[column sep=small]
		& 2 \arrow[rd, "h"] & \\
		1 \arrow[ru, "g"] \arrow[rr, "hg"'] & & 3
	\end{tikzcd} \quad \begin{tikzcd}[column sep=small]
		& 2 \arrow[rd, "h"] & \\
		0 \arrow[ru, "gf"] \arrow[rr, "h(gf)"'] & & 3.
	\end{tikzcd}\]
	This horn has a unique extension $\sigma:\Delta^3\to X$ such that
	$d^2(\sigma)\in X_2$ has the form
	% segment-id: o014.li2.chapter8.nerves-of-categories.q011
	\[\begin{tikzcd}[column sep=small]
		& 1 \arrow[rd, "hg"] & \\
		0 \arrow[ru, "f"] \arrow[rr, "h(gf)"'] & & 3.
	\end{tikzcd} \]
	By construction, however, this $2$-simplex also determines the composite
	of $f$ and $hg$, proving $(hg)f=h(gf)$.

	Thus $\mathcal{C}$ is a category. There is a canonical morphism
	$X\to\mathrm{N}\mathcal{C}$: from a given $\Delta^n\to X$, pull back
	along each $\Delta^1\xrightarrow{\{j,j+1\}}\Delta^n$ to obtain a chain of
	morphisms. By construction, $X_n\to\mathrm{N}\mathcal{C}_n$ is clearly a
	bijection for $n=0,1$. If $n\geq2$, choose $0<i<n$ and consider the
	commutative diagram
	% segment-id: o014.li2.chapter8.nerves-of-categories.q012
	\[\begin{tikzcd}
		\Hom(\Delta^n, X) \arrow[d] \arrow[r] & \Hom(\Delta^n, \mathrm{N}\mathcal{C}) \arrow[d] \\
		\Hom(\Lambda^n_i, X) \arrow[r] & \Hom(\Lambda^n_i, \mathrm{N}\mathcal{C}).
	\end{tikzcd}\]
	By the assumption and the unique inner horn-filling property, both vertical
	arrows are bijections. Moreover, $\Lambda^n_i$ can be expressed as the
	$\varinjlim$ of a family of $\Delta^{n-1}$ and $\Delta^{n-2}$ (that is,
	the gluing of $n$ faces), so induction shows that the arrow in the second
	row is also a bijection. This proves the result.
\end{proof}

% segment-id: o014.li2.chapter8.nerves-of-categories.p003
Notice also that in the argument (ii) $\implies$ (i), if there is a small
category $\mathcal{C}_0$ with $X=\mathrm{N}(\mathcal{C}_0)$, then the
category $\mathcal{C}$ constructed in the proof is exactly
$\mathcal{C}_0$.

% segment-id: o014.li2.chapter8.nerves-of-categories.p004
Write $\cate{Cat}$ for the category of all small categories, with functors as
morphisms. Every functor $F:\mathcal{C}\to\mathcal{C}'$ induces a morphism
$\mathrm{N}F:\mathrm{N}\mathcal{C}\to\mathrm{N}\mathcal{C}'$ in
$\cate{sSet}$, mapping $(f_1,\ldots,f_n)$ to
$(Ff_1,\ldots,Ff_n)$. This gives the nerve functor
$\mathrm{N}:\cate{Cat}\to\cate{sSet}$.

% segment-id: o014.li2.chapter8.nerves-of-categories.pr002
\begin{proposition}\label{prop:nerve-ff}
	The nerve functor $\mathrm{N}:\cate{Cat}\to\cate{sSet}$ is fully
	faithful.
\end{proposition}
% segment-id: o014.li2.chapter8.nerves-of-categories.pf002
\begin{proof}
	The proof of Proposition~\ref{prop:inner-horn-filling} has already shown
	how to reconstruct morphisms and their composition from the nerve of a
	category. Hence a morphism
	$\mathrm{N}(\mathcal{C})\to\mathrm{N}(\mathcal{C}')$ naturally induces a
	functor $\mathcal{C}\to\mathcal{C}'$. It is easy to see that this
	construction and the construction induced by $\mathrm{N}$ are mutually
	inverse.
\end{proof}

% segment-id: o014.li2.chapter8.nerves-of-categories.eg002
\begin{example}[Classifying space]\label{eg:classifying-space}
	\index{classifying space}
	\index[sym1]{BGamma@$\mathrm{B}\Gamma$}
	\index[sym1]{EGamma@$\mathrm{E}\Gamma$}
	Let $\Gamma$ be a monoid. Define the categories $\mathcal{B}\Gamma$ and
	$\mathcal{E}\Gamma$ as follows:
	% segment-id: o014.li2.chapter8.nerves-of-categories.q013
	\begin{center}\begin{tabular}{|c|c|c|c|}\hline
		Category & Object set & Morphisms & Composition of morphisms \\ \hline
		$\mathcal{B}\Gamma$ & $\{\star\}$ & $\End_{\mathcal{B}\Gamma}(\star)=\Gamma$ & multiplication in $\Gamma$ \\ \hline
		$\mathcal{E}\Gamma$ & $\Gamma$ & $\Hom_{\mathcal{E}\Gamma}(g,g')=\{h\in\Gamma:hg=g'\}$ & multiplication in $\Gamma$ \\ \hline
	\end{tabular}\end{center}

	Define the nerves
	$\mathrm{B}\Gamma:=\mathrm{N}((\mathcal{B}\Gamma)^{\opp})$ and
	$\mathrm{E}\Gamma:=\mathrm{N}((\mathcal{E}\Gamma)^{\opp})$.
	The simplicial set $\mathrm{B}\Gamma$ is called the classifying space of
	$\Gamma$. Taking $(\cdots)^{\opp}$ amounts to reversing the arrows in the
	definition of the nerve without changing their labels. Thus
	% segment-id: o014.li2.chapter8.nerves-of-categories.q014
	\begin{equation*}
		\begin{array}{|c|c|c|} \hline
		\text{Set} & \text{Element} & \text{Corresponding chain of morphisms in }\mathcal{B}\Gamma\text{ or }\mathcal{E}\Gamma \\ \hline
		(\mathrm{B}\Gamma)_n = \Gamma^n & (g_1, \ldots, g_n) & \star \xleftarrow{g_1} \cdots \leftarrow \star \xleftarrow{g_n} \star \\ \hline
		(\mathrm{E}\Gamma)_n = \Gamma^{n+1} & (g_1, \ldots, g_{n+1}) & g_1 \cdots g_{n+1} \xleftarrow{g_1} \cdots \leftarrow g_n g_{n+1} \xleftarrow{g_n} g_{n+1} \\ \hline
		\end{array}
	\end{equation*}
	Substitution into the definitions gives, for every $n\geq0$,
	% segment-id: o014.li2.chapter8.nerves-of-categories.q015
	\begin{equation*}
		\mathrm{B}\Gamma: \qquad \begin{aligned}
			d_i(g_1, \ldots, g_n) & =
			\begin{cases}
				(g_2, \ldots, g_n), & i = 0 \\
				(\ldots, g_i g_{i+1} , \ldots), & 0 < i < n \\
				(g_1, \ldots, g_{n-1}), & i = n,
			\end{cases} \\
			s_j(g_1, \ldots, g_n) & =
			\begin{cases}
				(1, g_1, \ldots, g_n), & j = 0 \\
				(\ldots, g_j, 1, g_{j+1}, \ldots), & 0 < j < n \\
				(g_1, \ldots, g_n, 1), & j = n,
			\end{cases}
		\end{aligned}
	\end{equation*}
	and
	% segment-id: o014.li2.chapter8.nerves-of-categories.q016
	\begin{equation*}
		\mathrm{E}\Gamma: \qquad \begin{aligned}
			d_i(g_1, \ldots, g_{n+1}) & =
			\begin{cases}
				(g_2, \ldots, g_{n+1}), & i = 0 \\
				(g_1, \ldots, g_i g_{i+1} , \ldots), & 0 < i \leq n,
			\end{cases} \\
			s_j(g_1, \ldots, g_{n+1}) & =
			\begin{cases}
				(1, g_1, \ldots, g_{n+1}), & j = 0 \\
				(\ldots, g_j, 1, g_{j+1}, \ldots), & 0 < j \leq n.
			\end{cases}
		\end{aligned}
	\end{equation*}

	There is a natural functor
	$(\mathcal{E}\Gamma)^{\opp}\to(\mathcal{B}\Gamma)^{\opp}$ mapping all
	objects to $\star$ and every morphism $h\in\Gamma$ to $h$. The induced
	morphism $\mathrm{E}\Gamma\to\mathrm{B}\Gamma$ simply maps
	$(g_1,\ldots,g_{n+1})$ to $(g_1,\ldots,g_n)$. The right multiplication
	action of $\Gamma$ on $(\mathrm{E}\Gamma)_n$ acts along its fibers, or in
	other words changes the datum $g_{n+1}$:
	% segment-id: o014.li2.chapter8.nerves-of-categories.q017
	\[ (g_1, \ldots, g_n, g_{n+1})g = (g_1, \ldots, g_n, g_{n+1}g). \]
	These actions clearly commute with $d_i$ and $s_j$. If $\Gamma$ is a
	group, the action is also free.

	In topology, $\mathrm{B}\Gamma$ (or its geometric realization
	$|\mathrm{B}\Gamma|$; see \S\ref{sec:geom-realization}) serves to
	classify $\Gamma$-torsors over spaces, while
	$\mathrm{E}\Gamma\to\mathrm{B}\Gamma$ gives the universal torsor.
	Example~\ref{eg:classifying-contractible} will show that
	$\mathrm{E}\Gamma$ is ``right-contractible.''
\end{example}
